\ifmostrarGabarito
\begin{minted}{python}
import pandas as pd

def relatorio_consumo_patrulha(df_missoes: pd.DataFrame) -> pd.DataFrame:
    """
    Agrupa missões por classe e zona de patrulha, calculando métricas e consumo por hora.
    """
    # Agrupamento multi-chave com agregações nomeadas via .agg()
    df_resumo = df_missoes.groupby(["classe_navio", "zona_patrulha"]).agg(
        horas_totais=("horas_navegadas", "sum"),
        horas_medias=("horas_navegadas", "mean"),
        consumo_total=("consumo_toneladas", "sum"),
        consumo_maximo=("consumo_toneladas", "max")
    ).reset_index()
    
    # Criação de coluna calculada de eficiência energética
    df_resumo["taxa_consumo_hora"] = df_resumo["consumo_total"] / df_resumo["horas_totais"]
    
    # Arredondamento numérico uniforme para 2 casas decimais
    df_resumo = df_resumo.round(2)
    
    return df_resumo
\end{minted}

\textbf{Explicação:}
A técnica de agregações nomeadas (\emph{Named Aggregation}) no método \texttt{.agg(novo\_nome=('coluna', 'operacao'))} gera um DataFrame plano e autoexplicativo, eliminando índices hierárquicos complexos. A taxa derivada por hora sintetiza a métrica operacional necessária para o planejamento logístico da força naval.
\fi
